% Anti-Tech Presentation
\documentclass[11pt]{beamer}
\usetheme{metropolis}


% -----------------------
% Packages
% -----------------------
\usepackage{graphicx}
\usepackage{amsmath, amssymb}
\usepackage{booktabs}
\usepackage{hyperref}
\usepackage{algorithm}
\usepackage{algpseudocode}
\usepackage{listings}
\lstset{
  basicstyle=\ttfamily\small,
  columns=fullflexible,
  frame=single,
  breaklines=true
}


% -----------------------
% Title info
% -----------------------
\title{Assessing Architecture Choices Using Hessian Eigndensity}
\date{\today}
\author{Flavio D. Gualtieri}
\begin{document}


% =======================
% Slide 1: Title
% =======================
\begin{frame}
  \titlepage
\end{frame}


% =======================
% Slide 2: Motivation / Aim (optional but helpful)
% =======================
\begin{frame}{Goal}
    \begin{itemize}
      \item Objective: test whether tracking the \textbf{evolution of Hessian eigendensity} during training
      can provide insight into \textbf{architecture/activation choices}.
      \item Proof-of-concept: small MNIST network, compare activation functions.
      \item Output: eigendensity spectra snapshots throughout training.
    \end{itemize}
\end{frame}


% =======================
% Slide 3: General pseudocode
% =======================
\begin{frame}{Estimating Eigendensity (Pseudocode)}
    \begin{algorithmic}
        \State Forward pass; compute loss $\ell(\theta)$
        \State Generate $k$ IID probe vectors $\{v_1, ..., v_k\}$
        \For{$i \in [1, \dots, k]$}
            \State Run $m$ - step Lanczos on $H$ with $v_i$
            \State Obtain $T_m \in \mathbb{R}^{m\times m}$
            \State Compute $T_m = ULU^\top$
            \For{$r \in [1,\dots,m]$}
                \State $w_r \gets (U_{1r})^2$
                \State $\lambda_r \gets (L_{rr})^m_{r=1}$
            \EndFor
            \State $\hat{\phi}^{v_i}(t) \gets \sum_{r=1}^m w_r f(\lambda_r;t,\sigma^2)$
        \EndFor
        \State $\hat{\phi}(t) \gets \frac{1}{k}\sum_{i=1}^k \hat{\phi}_i(t)$
    \end{algorithmic}

    \vspace{0.2em}
    \footnotesize
    Note: $f(\lambda; t, \sigma^2) = \frac{1}{\sigma\sqrt{2\pi}}e^{-\frac{(t-\lambda)^2}{2\sigma^2}}$
\end{frame}


% =======================
% Slide 4: Experimental setup
% =======================
\begin{frame}{Experimental Setup}
    \begin{columns}[T,onlytextwidth]
      \column{0.52\textwidth}
      \begin{block}{Model / Data}
      \begin{itemize}
        \item Dataset: MNIST
        \item Architecture: \([784,\,128,\,10]\)
        \item Activations: \(\{\)sigmoid, ReLU, tanh\(\}\)
      \end{itemize}
      \end{block}
    
      \column{0.48\textwidth}
      \begin{block}{Training / Logging}
      \begin{itemize}
        \item Optimizer: Adam
        \item Learning rate: \(10^{-3}\)
        \item Epochs: 50
        \item Eigdensity snapshot: every 10 epochs
      \end{itemize}
      \end{block}
    \end{columns}
    
    \vspace{0.4em}
    \footnotesize
    Epoch budget picked to be as small as possible (within reason) after checking convergence trends, avoiding needless compute.
\end{frame}


% =======================
% Slide 5: Procedure timeline (optional)
% =======================
\begin{frame}{Procedure Summary}
    \begin{enumerate}
      \item Choose activation \(\in\{\)sigmoid, ReLU, tanh\(\}\)
      \item Train for 50 epochs (Adam, lr \(=10^{-3}\))
      \item At epochs \(\{10,20,30,40,50\}\): estimate Hessian eigendensity
      \item Plot spectra and compare across activations
    \end{enumerate}
    
    \vspace{0.5em}
    \begin{block}{Result}
    A set of eigendensity spectra over training time, enabling qualitative comparison across choices.
    \end{block}
\end{frame}


% =======================
% Slide 6: Sigmoid spectra
% =======================
\begin{frame}{Eigenvalue Density Spectra — Sigmoid}
    \centering
    \begin{tabular}{ccc}
        \includegraphics[width=0.3\textwidth]{sigmoid_phi_epoch_1.png} &
        \includegraphics[width=0.3\textwidth]{sigmoid_phi_epoch_10.png} &
        \includegraphics[width=0.3\textwidth]{sigmoid_phi_epoch_20.png} \\
        Epoch 1 & Epoch 10 & Epoch 20 \\
        \vspace{0.2cm} \\
        \includegraphics[width=0.3\textwidth]{sigmoid_phi_epoch_30.png} &
        \includegraphics[width=0.3\textwidth]{sigmoid_phi_epoch_40.png} &
        \includegraphics[width=0.3\textwidth]{sigmoid_phi_epoch_50.png} \\
        Epoch 30 & Epoch 40 & Epoch 50
    \end{tabular}
    
    \vspace{0.3cm}
    {\footnotesize Note: test accuracy: 0.9789}
\end{frame}


% =======================
% Slide 7: ReLU spectra
% =======================
\begin{frame}{Eigenvalue Density Spectra — ReLU}
    \centering
    \begin{tabular}{ccc}
        \includegraphics[width=0.3\textwidth]{relu_phi_epoch_1.png} &
        \includegraphics[width=0.3\textwidth]{relu_phi_epoch_10.png} &
        \includegraphics[width=0.3\textwidth]{relu_phi_epoch_20.png} \\
        Epoch 1 & Epoch 10 & Epoch 20 \\
        \vspace{0.2cm} \\
        \includegraphics[width=0.3\textwidth]{relu_phi_epoch_30.png} &
        \includegraphics[width=0.3\textwidth]{relu_phi_epoch_40.png} &
        \includegraphics[width=0.3\textwidth]{relu_phi_epoch_50.png} \\
        Epoch 30 & Epoch 40 & Epoch 50
    \end{tabular}

    \vspace{0.3cm}
    {\footnotesize Note: test accuracy: 0.9742}
\end{frame}


% =======================
% Slide 8: Tanh spectra
% =======================
\begin{frame}{Eigenvalue Density Spectra — Tanh}
    \centering
    \begin{tabular}{ccc}
        \includegraphics[width=0.3\textwidth]{tanh_phi_epoch_1.png} &
        \includegraphics[width=0.3\textwidth]{tanh_phi_epoch_10.png} &
        \includegraphics[width=0.3\textwidth]{tanh_phi_epoch_20.png} \\
        Epoch 1 & Epoch 10 & Epoch 20 \\
        \vspace{0.2cm} \\
        \includegraphics[width=0.3\textwidth]{tanh_phi_epoch_30.png} &
        \includegraphics[width=0.3\textwidth]{tanh_phi_epoch_40.png} &
        \includegraphics[width=0.3\textwidth]{tanh_phi_epoch_50.png} \\
        Epoch 30 & Epoch 40 & Epoch 50
    \end{tabular}

    \vspace{0.3cm}
    {\footnotesize Note: test accuracy: 0.9753}
\end{frame}


% =======================
% Slide 9: Spectral Ratio Evolution
% =======================
\begin{frame}{Spectral Ratio Evolution Across Activations}
    \centering
    \begin{tabular}{ccc}
        \includegraphics[width=0.3\textwidth]{sigmoid_kappa_vs_epoch.png} &
        \includegraphics[width=0.3\textwidth]{relu_kappa_vs_epoch.png} &
        \includegraphics[width=0.3\textwidth]{tanh_kappa_vs_epoch.png} \\
        Sigmoid & ReLU & Tanh
    \end{tabular}

    \vspace{0.3cm}
    {\footnotesize Evolution during training of 
    $\lambda_{\max} / \lambda_{\min}$, where $\lambda_{\max}$ is the largest Hessian eigenvalue and 
    $\lambda_{\min}$ is the smallest positive non-zero eigenvalue.}
\end{frame}


\end{document}
